\documentclass{beamer}
\usepackage{booktabs,tikz}
\usetheme{Berlin}
\title{Risk and Return}
\begin{document}
\maketitle

\section{Introduction}
\begin{frame}{Risk and Return}
	\begin{itemize}
		\item Investors prefer more return and less risk
		\item In this lecture, we will focus on two concepts related to investors' preferences over return and risk
		\begin{enumerate}
			\item The risk premium, which is the return we expect to receive for taking a certain amount of risk
			\item The Sharpe ratio, which tells us how much return we receive for each unit of risk we take
		\end{enumerate}
	\end{itemize}
\end{frame}

\section{Risk Aversion}

\begin{frame}{Risk Aversion}
	\begin{center}
		\renewcommand{\arraystretch}{1.5}
		\begin{tabular}{cc|cc}
			\multicolumn{2}{c}{\textbf{Bet A}} & \multicolumn{2}{c}{\textbf{Bet B}} \\
			Amount & Probability & Amount & Probability \\ \hline
			-10k USD & 0.5 & 10k USD & 1. \\
			30k USD & 0.5 & & \\ \hline
			\multicolumn{2}{c}{Average = 10k USD} & \multicolumn{2}{c}{Average = 10k}
		\end{tabular}
	\end{center}
\begin{center}
A risk averse investor prefers \textbf{Bet B} to \textbf{Bet A}. While both bets return 10k USD on average, \textbf{Bet A} is significantly riskier.
\end{center}
\end{frame}

\begin{frame}{Risk Aversion}
	\begin{center}
		\renewcommand{\arraystretch}{1.5}
		\begin{tabular}{cc|cc}
			\multicolumn{2}{c}{\textbf{Bet A}} & \multicolumn{2}{c}{\textbf{Bet B}} \\
			Amount & Probability & Amount & Probability \\ \hline
			-10k USD & 0.5 & 10k USD & 1. \\
			40k USD & 0.5 & & \\ \hline
			\multicolumn{2}{c}{Average = 15k USD} & \multicolumn{2}{c}{Average = 10k USD}
		\end{tabular}
	\end{center}
\begin{center}
She will only accept \textbf{Bet A} to \textbf{Bet B} if she is compensated for doing so, in this example, by earning an additional 5k USD.
\end{center}
\end{frame}

\section{Risk Premia}

\begin{frame}{Risk-Free Asset}
	\begin{itemize}
		\item The risk-free rate is the rate of return that can be earned with certainty
		\item The risk-free asset is the asset that yields a risk-free return
		\item We take US Treasuries to be risk-free assets and Treasury yields to be the risk-free return
		\item Let's call the risk-free rate $R_{F}$ 
	\end{itemize}
\end{frame}

\begin{frame}{Risk Premia}
	\begin{itemize}
		\item The risk premium is the expected excess return over the risk-free rate
		\begin{equation}
			\text{RP}[R_{P}]=\text{E}[R_{P}]-R_{F}
		\end{equation}
		\item It is the return we earn for taking risk: higher risk stocks will generally have higher risk premia, safer stocks will generally have lower risk premia
		\item If the portfolio $P$ is risk-free, its risk-premium will be zero
	\end{itemize}
\end{frame}

\begin{frame}{Risk Premia}
	The mean monthly return on the market from 1929 to 2019 was 0.9379\%. The mean monthly return on US T-bills was 0.2732\%. What was the market risk premium over this period?\\
	\pause
	{\color{blue}
		The monthly market risk premium was
		\begin{equation}
			\text{RP}[R_{\text{Market}}]=0.9379\%-0.2732\%=0.6647\%
		\end{equation}
		This is how much the market had to compensate investors for the risk they were taking.
	}
\end{frame}

\section{Sharpe Ratio}

\begin{frame}{Sharpe Ratios}
	\begin{itemize}
		\item The Sharpe ratio is the ratio of a portfolio's risk premium to its standard deviation:
			\begin{equation}
				\text{SR}[R_{P}]=\frac{\text{RP}[R_{P}]}{\text{SD}[R_{P}]}=\frac{\text{E}[R_{P}]-R_{F}}{\text{SD}[R_{P}]}
			\end{equation}
			where $\text{E}[R_{P}]$ and $\text{SD}[R_{P}]$ are the mean and standard deviation of the return on the portfolio respectively 
		\item It is the return we earn for each unit of risk we take
	\end{itemize}
\end{frame}

\begin{frame}{Sharpe Ratios}
	The mean monthly standard deviation of market returns from 1929 to 2019 was 5.3324\%. What was the market Sharpe ratio over this period?\\
	\pause
	{\color{blue}
		The monthly market risk premium was
		\begin{equation}
			\text{SR}[R_{\text{Market}}]=\frac{.6647\%}{5.3324\%}=0.1247
		\end{equation}
}
\end{frame}

\end{document}
